\chapter{Discussion, limites et recommandations}

\section{Synthèse des résultats}

Ce travail a permis de construire, à partir de neuf dossiers de crédit réels du Centre d'Affaires
Al Istiqlal, un système hybride combinant une grille de scoring experte (Chapitre~4), un modèle
d'apprentissage automatique évalué en LOOCV (Chapitre~5), et un moteur de décision logique
explicite (Chapitre~6). Trois résultats principaux se dégagent :

\begin{enumerate}
    \item Le système hybride reproduit correctement la notation MNS2 réelle à un cran près pour
    huit dossiers sur neuf, et isole automatiquement le seul cas de divergence marquée (Feuil2)
    pour investigation qualitative — validant la capacité du dispositif à \emph{signaler} les cas
    atypiques plutôt qu'à simplement produire un score ;
    \item L'analyse d'importance des variables (Random Forest) fait ressortir des facteurs
    financiers cohérents avec la théorie du risque de crédit — structure de l'actif immobilisé et
    liquidité immédiate en tête — apportant un premier niveau d'explicabilité globale au système ;
    \item Le moteur de décision, en s'abstenant de trancher automatiquement lorsque la confiance
    du modèle est insuffisante (zone « à étudier »), traduit concrètement l'articulation entre
    automatisation et jugement humain qui structure la problématique de ce projet (Chapitre~1).
\end{enumerate}

\section{Limites méthodologiques}

Trois limites, déjà évoquées ponctuellement dans les chapitres précédents, méritent d'être
rassemblées et hiérarchisées ici.

\begin{limite}[1 — Taille de l'échantillon]
Avec 9 dossiers (7 labellisés pour la détection de défaut), toute métrique de performance
statistique doit être interprétée comme indicative. Le LOOCV, bien qu'étant le protocole le plus
rigoureux disponible à cette échelle, ne peut se substituer à une validation sur un échantillon
de taille suffisante pour estimer une performance généralisable.
\end{limite}

\begin{limite}[2 — Variable cible proxy]
En l'absence d'historique réel d'incidents de paiement, la cible de défaut a été construite par
une coupure statistique (K-means) sur le score MNS2 existant. Cette approche introduit une
circularité partielle : le modèle apprend, dans une certaine mesure, à retrouver une structure
déjà présente dans le score qu'il est censé compléter, plutôt qu'à prédire un événement de
défaut indépendant.
\end{limite}

\begin{limite}[3 — Paramètres non calibrés]
Les pondérations de la grille experte (Chapitre~4) et les seuils de décision (60\,\%/25\,\%,
Chapitre~6) sont des hypothèses de travail raisonnées mais non calibrées sur des données de
validation indépendantes, faute d'échantillon suffisant pour réserver un sous-ensemble à cet
effet.
\end{limite}

\section{Recommandations pour une industrialisation}

Trois recommandations découlent directement des limites identifiées :

\begin{itemize}
    \item \textbf{Constitution d'un échantillon élargi} : le workflow applicatif (Chapitre~6,
    Section 6.4) a été conçu pour capitaliser sur les nouveaux dossiers traités par les analystes.
    Chaque dossier supplémentaire, une fois validé par la Direction, peut enrichir l'échantillon
    d'entraînement et permettre, à terme, un passage d'une validation LOOCV à une validation
    train/test classique, plus robuste ;
    \item \textbf{Remplacement de la cible proxy} : dès qu'un historique réel d'incidents de
    paiement devient accessible (même partiellement), il devrait se substituer à la coupure
    K-means actuelle pour fonder l'apprentissage sur un événement de défaut réel plutôt que sur
    une structure interne au score existant ;
    \item \textbf{Calibration collaborative de la grille experte} : une fois un échantillon plus
    large disponible, les pondérations pourraient être ajustées de façon supervisée (par exemple,
    par une régression contrainte visant à minimiser l'écart entre score expert et notation MNS2
    réelle), tout en conservant l'intelligibilité de la grille.
\end{itemize}

\section{Conclusion générale}

Ce projet répond à la problématique posée au Chapitre~1 — reconstruire et expliquer une
notation bancaire perçue comme une boîte noire — par un système qui articule explicitement trois
niveaux : une logique de règles transparentes (grille experte), une logique statistique
(apprentissage automatique et importance des variables), et une couche de décision qui assume
ses limites en renvoyant les cas incertains à l'arbitrage humain plutôt que de forcer un verdict.
Les résultats obtenus, bien que devant être lus avec la prudence qu'impose la taille de
l'échantillon, démontrent la faisabilité méthodologique de la démarche et fournissent une base
concrète — code, base de données anonymisée, application fonctionnelle — sur laquelle une montée
en échelle pourra s'appuyer.
